%% Generated by Sphinx.
\def\sphinxdocclass{report}
\documentclass[letterpaper,10pt,brazil]{sphinxmanual}
\ifdefined\pdfpxdimen
   \let\sphinxpxdimen\pdfpxdimen\else\newdimen\sphinxpxdimen
\fi \sphinxpxdimen=.75bp\relax
\ifdefined\pdfimageresolution
    \pdfimageresolution= \numexpr \dimexpr1in\relax/\sphinxpxdimen\relax
\fi
%% let collapsible pdf bookmarks panel have high depth per default
\PassOptionsToPackage{bookmarksdepth=5}{hyperref}

\PassOptionsToPackage{booktabs}{sphinx}
\PassOptionsToPackage{colorrows}{sphinx}

\PassOptionsToPackage{warn}{textcomp}
\usepackage[utf8]{inputenc}
\ifdefined\DeclareUnicodeCharacter
% support both utf8 and utf8x syntaxes
  \ifdefined\DeclareUnicodeCharacterAsOptional
    \def\sphinxDUC#1{\DeclareUnicodeCharacter{"#1}}
  \else
    \let\sphinxDUC\DeclareUnicodeCharacter
  \fi
  \sphinxDUC{00A0}{\nobreakspace}
  \sphinxDUC{2500}{\sphinxunichar{2500}}
  \sphinxDUC{2502}{\sphinxunichar{2502}}
  \sphinxDUC{2514}{\sphinxunichar{2514}}
  \sphinxDUC{251C}{\sphinxunichar{251C}}
  \sphinxDUC{2572}{\textbackslash}
\fi
\usepackage{cmap}
\usepackage[T1]{fontenc}
\usepackage{amsmath,amssymb,amstext}
\usepackage{babel}



\usepackage{tgtermes}
\usepackage{tgheros}
\renewcommand{\ttdefault}{txtt}



\usepackage[Sonny]{fncychap}
\ChNameVar{\Large\normalfont\sffamily}
\ChTitleVar{\Large\normalfont\sffamily}
\usepackage{sphinx}

\fvset{fontsize=auto}
\usepackage{geometry}


% Include hyperref last.
\usepackage{hyperref}
% Fix anchor placement for figures with captions.
\usepackage{hypcap}% it must be loaded after hyperref.
% Set up styles of URL: it should be placed after hyperref.
\urlstyle{same}

\addto\captionsbrazil{\renewcommand{\contentsname}{Contents:}}

\usepackage{sphinxmessages}
\setcounter{tocdepth}{1}



\title{Lógica de programação}
\date{14 ago. 2025}
\release{1.0}
\author{Hamilton Ventura Rodrigues}
\newcommand{\sphinxlogo}{\vbox{}}
\renewcommand{\releasename}{Release}
\makeindex
\begin{document}

\ifdefined\shorthandoff
  \ifnum\catcode`\=\string=\active\shorthandoff{=}\fi
  \ifnum\catcode`\"=\active\shorthandoff{"}\fi
\fi

\pagestyle{empty}
\sphinxmaketitle
\pagestyle{plain}
\sphinxtableofcontents
\pagestyle{normal}
\phantomsection\label{\detokenize{index::doc}}

\subsubsection*{Introdução}

\sphinxAtStartPar
Bem\sphinxhyphen{}vindo(a) à \sphinxstyleemphasis{introdução à lógica de programação}. Aqui reúno um compêndio de exercícios e
explicações com ferramentas conhecidas que podem facilitar a criança, adolescente ou adulto na
introdução ao mundo da programação. Possuímos atualmente a versão web e PDF, caso tenha interesse.

\sphinxAtStartPar
Como está dividido esse documento?
\begin{itemize}
\item {} 
\sphinxAtStartPar
Navegação ao Sistema Operacional

\item {} 
\sphinxAtStartPar
Comandos CLI no Sistema Operacional

\item {} 
\sphinxAtStartPar
REGEX

\end{itemize}
\subsubsection*{Engine Defold}

\sphinxAtStartPar
Engine para desenvolvimento de jogos, focada em desenvolvimento de jogos de forma visual.
\subsubsection*{Carga horária}

\sphinxAtStartPar
As aulas de Programação com jogos totalizam 20 horas, sendo uma hora e meia a cada
encontro, duas vezes durante a semana, totalizando 13 aulas em 2 meses.
\subsubsection*{Objetivos Gerais}

\sphinxAtStartPar
O curso de programação com games busca desenvolver nos alunos o pensamento
computacional e o raciocínio lógico, de maneira a resolver problemas mais complexos de
forma lúdica.
\subsubsection*{Objetivos Específicos}

\sphinxAtStartPar
Operações básicas no Sistema Operacional:
Aprender a utilizar o básico do sistema operacional, ajudando a introduzir os alunos à
informática básica, promovendo um melhor aprendizado nas ferramentas propostas
no curso (que exigem conhecimento básico).

\sphinxAtStartPar
Aprender a desenvolver soluções na linguagem lua:
Os princípios essenciais a serem aprendidos: condições, laços, variáveis,
acumuladores, eventos, transmissão de informação, edição de imagens dos
personagens. Utilização livre e espontânea da ferramenta.

\sphinxAtStartPar
Como princípios essenciais a serem desenvolvidos pelos alunos está, em primeiro
lugar, entender o funcionamento básico da ferramenta: navegação em tela, salvar e
abrir arquivos, criação de pastas no Windows, salvar imagens do Google no
computador pessoal. As ações poderão ser desenvolvidas utilizando como base os
princípios do pensamento lógico. No decorrer das aulas o aluno aprenderá a criar
personagens, cenários, mudança de cenário, disparos de projéteis, criação de
barras de vida e contadores de moedas. Durante as aulas os alunos aprenderão a
ativar animações em determinadas situações, trabalhando com eventos.

\sphinxstepscope


\chapter{Introdução ao Curso}
\label{\detokenize{modelo_aula_02:introducao-ao-curso}}\label{\detokenize{modelo_aula_02::doc}}

\section{Introdução a lógica de programação com jogos}
\label{\detokenize{modelo_aula_02:introducao-a-logica-de-programacao-com-jogos}}
\sphinxAtStartPar
\sphinxstylestrong{Tempo de Aula:} 1h30min
\subsubsection*{Recursos Didáticos}
\begin{itemize}
\item {} 
\sphinxAtStartPar
Projetor: Exibição de exemplos, Slides e exercícios

\item {} 
\sphinxAtStartPar
Quadro branco: Exibição de mecânicas de forma desenhada

\item {} 
\sphinxAtStartPar
Computadores: Computadores com os softwares citados a baixo instaldos e internet.

\end{itemize}
\subsubsection*{Metodologia}
\begin{itemize}
\item {} 
\sphinxAtStartPar
Expositiva, aplicação de exercícios e solução de problemas.

\item {} 
\sphinxAtStartPar
Exercícios Guiados em programação.

\item {} 
\sphinxAtStartPar
Atividade de exploração Criativa.

\item {} 
\sphinxAtStartPar
Estimulo a participação.

\end{itemize}
\subsubsection*{Objetivos}
\begin{itemize}
\item {} 
\sphinxAtStartPar
Compreender conceitos básicos que estão por trás da criação de jogos.

\item {} 
\sphinxAtStartPar
Conhecer e ter uma visão Geral das ferramentas gráficas para criação de sprites e artes digitais.

\item {} 
\sphinxAtStartPar
Desenvolver raciocínio lógico e habilidades criativas por meio de exercícios práticos.

\end{itemize}
\subsubsection*{Conteúdo Programático}

\sphinxAtStartPar
\sphinxstylestrong{Ambientação}
\begin{itemize}
\item {} 
\sphinxAtStartPar
Ligar computador.

\item {} 
\sphinxAtStartPar
Pesquisar conteúdos no navegador web.

\item {} 
\sphinxAtStartPar
Criação de Pastas.

\end{itemize}

\sphinxAtStartPar
\sphinxstylestrong{Engine}
\begin{itemize}
\item {} 
\sphinxAtStartPar
Iniciar um projeto no Defold;

\item {} 
\sphinxAtStartPar
Navegar nas áreas: área de trabalho, Terminal.

\end{itemize}

\sphinxAtStartPar
{\color{red}\bfseries{}*}Para os alunos irem se ambientando e familiarizando com a ferramentas é necessário fazer um exploratório
principlamente do Defold.

\sphinxAtStartPar
\sphinxstylestrong{Fundamentos dos Jogos}
\begin{itemize}
\item {} 
\sphinxAtStartPar
Termos e definições de tipos de imagens (JPG,PNG,GIF E SVG)

\item {} 
\sphinxAtStartPar
Termos e definições: Sprites, Plataformas, Engines, Tiles

\item {} 
\sphinxAtStartPar
Exemplos de Jogos que foram desenvolvidos utilizando a linguagem.

\end{itemize}

\sphinxAtStartPar
\sphinxstylestrong{Ferramentas Gráficas}
\begin{itemize}
\item {} 
\sphinxAtStartPar
Libre Sprite

\item {} 
\sphinxAtStartPar
Inkscape

\end{itemize}
\subsubsection*{Descrição}

\sphinxAtStartPar
A aula abordará os conceitos introdutórios sobre jogos e seu funcionamento,
incluindo alguns exemplos. Diversos termos serão abordados na aula, introduzindo o aluno
ao mundo do desenvolvimento de jogos suas ferramentas principais.

\sphinxAtStartPar
Ao final, será reservado um momento para que as crianças possam explorar livremente essas ferramentas,
incentivando a criatividade, a experimentação e o desenvolvimento de habilidades artísticas e técnicas.
Após a apresentação dos Slides é necessário deixar pelo menos 30 minutos para as crianças brincarem.
\subsubsection*{Avaliação}


\section{Execução dos exercícios em sala de aula.}
\label{\detokenize{modelo_aula_02:execucao-dos-exercicios-em-sala-de-aula}}
\sphinxAtStartPar
\sphinxstylestrong{Exercício 1: Criar um exemplo de saída}
Crie um arquivo no vscode e digite um código que exiba Hellow Word na tela e seguido do seu nome.

\sphinxAtStartPar
\sphinxstylestrong{Exercício 2: Criar código direto no terminal}
Crie uma saida escrita Hellow Word mas diretamente na tela;

\sphinxAtStartPar
—
\subsubsection*{Bibliografia}

\sphinxAtStartPar
LUA. Lua.org. Rio de Janeiro: PUC\sphinxhyphen{}Rio, {[}s.d.{]}.
Disponível em: \sphinxurl{https://www.lua.org/}. Acesso em: 14 ago. 2025.

\sphinxAtStartPar
LIBRESPRITE. LibreSprite. {[}S. l.{]}: LibreSprite, {[}s.d.{]}.
Disponível em: \sphinxurl{https://libresprite.github.io/}. Acesso em: 14 ago. 2025.

\sphinxAtStartPar
INKSCAPE. Inkscape. {[}S. l.{]}: Inkscape, {[}s.d.{]}.
Disponível em: \sphinxurl{https://inkscape.org/}. Acesso em: 14 ago. 2025.

\sphinxAtStartPar
Rabin, S. (n.d.). Introdução ao desenvolvimento de games (Cengage Le, Vol. 1).
\sphinxurl{https://www.amazon.com.br}/Introdução\sphinxhyphen{}desenvolvimento\sphinxhyphen{}games\sphinxhyphen{}Entendendo\sphinxhyphen{}universo/dp/852211143X

\sphinxAtStartPar
Rita, S. (2017). Treinamento em Lógica de Programação (R. Sandra (ed.); Universo d).
\sphinxurl{https://br.99ebooks.net/download/treinamento-em-logica-de-programacao/}

\sphinxstepscope


\chapter{Scratch \sphinxhyphen{} introdução ao código}
\label{\detokenize{modelo_aula_03:scratch-introducao-ao-codigo}}\label{\detokenize{modelo_aula_03::doc}}
\sphinxAtStartPar
\sphinxstylestrong{Tempo de Aula:} 1h30min
\subsubsection*{Recursos Didáticos}

\sphinxAtStartPar
Projetor, quadro branco, computadores.
\subsubsection*{Metodologia}

\sphinxAtStartPar
Expositiva, aplicação de exercícios / solução de problemas.
\subsubsection*{Conteúdo Programático}

\sphinxAtStartPar
Como Aplicar movimentos aos personagens
APlicar o bloco de comando Repita.
Aplicar o bloco de comando SE
Utilização dos blocos de comandos sensores.
Utilizar o bloco de comando trocar fantasia
Inserir o cenário do jogo/animação
Salvando o projeto Projeto.
\subsubsection*{Descrição}

\sphinxAtStartPar
No primeiro momento deve\sphinxhyphen{}se explicar o conceito do Bloco de comando repita, fazendo com
que o personagem repita os movimentos até encostar na borda da tela.
Logo depois dentro do bloco de comando repita, utilizar o comando trocar fantasias.
Deve\sphinxhyphen{}se estimular os alunos a editar os personagens para simular animação, e ao mesmo
tempo modificar as fantasias para dar ilusão de movimento. Após essa parte ensinar a
Modificar o cenário de fundo com algum da escolha do aluno.
\subsubsection*{Avaliação}

\sphinxAtStartPar
Execução dos exercícios em sala de aula.
\subsubsection*{Exercícios de Desafio}

\sphinxAtStartPar
\sphinxstylestrong{Exercício 1: Criar uma animação simples}
Crie um projeto onde o personagem escolhido faça pelo menos dois movimentos diferentes (exemplo: andar para frente e girar).
Use blocos de código para controlar esses movimentos e troque as fantasias para simular animação.

\sphinxAtStartPar
\sphinxstylestrong{Exercício 2: Explorar a área de fantasias e sons}
Adicione uma nova fantasia ao personagem ou crie uma fantasia personalizada.
Depois, escolha ou grave um som e faça o personagem reproduzi\sphinxhyphen{}lo quando a bandeira verde for clicada.

\sphinxAtStartPar
\sphinxstylestrong{Exercício 3: Salvar e compartilhar o projeto}
Salve o projeto com um nome criativo.
Caso tenha acesso à internet, publique o projeto no site do Scratch e copie o link para compartilhar com um colega ou professor.
\subsubsection*{Bibliografia}

\sphinxAtStartPar
SCIRRA, L. (2022). Manual Construct.
\sphinxurl{https://www.construct.net/en/make-games/manuals/construct-3}

\sphinxAtStartPar
Rabin, S. (n.d.). Introdução ao desenvolvimento de games (Cengage Le, Vol. 1).
\sphinxurl{https://www.amazon.com.br}/Introdução\sphinxhyphen{}desenvolvimento\sphinxhyphen{}games\sphinxhyphen{}Entendendo\sphinxhyphen{}universo/dp/852211143X

\sphinxAtStartPar
Rita, S. (2017). Treinamento em Lógica de Programação (R. Sandra (ed.); Universo d).
\sphinxurl{https://br.99ebooks.net/download/treinamento-em-logica-de-programacao/}

\sphinxstepscope


\chapter{Scratch \sphinxhyphen{} Eventos de interação}
\label{\detokenize{modelo_aula_04:scratch-eventos-de-interacao}}\label{\detokenize{modelo_aula_04::doc}}
\sphinxAtStartPar
\sphinxstylestrong{Tempo de Aula:} 1h30min
\subsubsection*{Recursos Didáticos}

\sphinxAtStartPar
Projetor, quadro branco, computadores.
\subsubsection*{Metodologia}

\sphinxAtStartPar
Expositiva, aplicação de exercícios / solução de problemas.
\subsubsection*{Conteúdo Programático}
\begin{itemize}
\item {} 
\sphinxAtStartPar
Fazer dois ou mais personagens/objetos interagirem.

\item {} 
\sphinxAtStartPar
Programar teclas

\end{itemize}
\subsubsection*{Descrição}

\sphinxAtStartPar
Deve\sphinxhyphen{}se gastar pelo menos 20 a 30 minutos reforçando os conteúdos com alguns
exercícios. Após isso deve ser feito um exercício em conjunto para que o aluno entenda
como funciona o comando transmitir e aprenda a editar códigos em mais de um
personagem.
A programação das teclas do teclado podem ser feitas com o bloco de comando “Quando a
tecla for precionada”, podendo assim configurar os movimentos de esquerda, direita, cima e
baixo.
\subsubsection*{Avaliação}

\sphinxAtStartPar
Execução dos exercícios em sala de aula.
\subsubsection*{Exercícios de Desafio}

\sphinxAtStartPar
\sphinxstylestrong{Desafio 1: Criar uma conversa entre personagens}
Faça dois personagens enviarem mensagens entre si usando o bloco “Transmitir” e “Quando eu receber”.
Programe para que um personagem diga algo, e o outro responda após receber a mensagem.

\sphinxAtStartPar
\sphinxstylestrong{Desafio 2: Controle completo do personagem com teclado}
Programe um personagem para se movimentar em todas as direções (cima, baixo, esquerda, direita) usando as teclas correspondentes.
Adicione efeitos visuais, como troca de fantasia ou sons, para cada movimento realizado.
\subsubsection*{Bibliografia}

\sphinxAtStartPar
SCIRRA, L. (2022). Manual Construct.
\sphinxurl{https://www.construct.net/en/make-games/manuals/construct-3}

\sphinxAtStartPar
Rabin, S. (n.d.). Introdução ao desenvolvimento de games (Cengage Le, Vol. 1).
\sphinxurl{https://www.amazon.com.br}/Introdução\sphinxhyphen{}desenvolvimento\sphinxhyphen{}games\sphinxhyphen{}Entendendo\sphinxhyphen{}universo/dp/852211143X

\sphinxAtStartPar
Rita, S. (2017). Treinamento em Lógica de Programação (R. Sandra (ed.); Universo d).
\sphinxurl{https://br.99ebooks.net/download/treinamento-em-logica-de-programacao/}

\sphinxstepscope


\chapter{Introdução ao Scratch: Variáveis}
\label{\detokenize{modelo_aula_05:introducao-ao-scratch-variaveis}}\label{\detokenize{modelo_aula_05::doc}}
\sphinxAtStartPar
\sphinxstylestrong{Tempo de Aula:} 1h30min
\subsubsection*{Recursos Didáticos}

\sphinxAtStartPar
Projetor, quadro branco, computadores.
\subsubsection*{Metodologia}

\sphinxAtStartPar
Expositiva, aplicação de exercícios / solução de problemas.
\subsubsection*{Conteúdo Programático}

\sphinxAtStartPar
Utilizar as estruturas básicas do Scratch para um desafio final utilizando variáveis
para contar a quantidade de moedas recolhidas pelo gatinho.
\subsubsection*{Descrição}

\sphinxAtStartPar
Para esse desafio final, alguns elementos são necessários para o aluno conseguir entender
a lógica do construct:
1 \sphinxhyphen{} Condições (Bloco SE)
2 \sphinxhyphen{} Comparar Variáveis
Para tanto será feita uma revisão para posteriormente ser lançado o desafio.
\subsubsection*{Avaliação}

\sphinxAtStartPar
Execução dos exercícios em sala de aula.
\subsubsection*{Exercícios de Desafio}

\sphinxAtStartPar
\sphinxstylestrong{Desafio 1: Criar uma condição simples}
Utilize o bloco SE para criar um programa onde, se uma variável atingir um valor específico,
uma ação seja executada (exemplo: mudar a cor de um objeto ou mostrar uma mensagem).

\sphinxAtStartPar
\sphinxstylestrong{Desafio 2: Comparação de variáveis}
Programe duas variáveis para serem comparadas.
Quando as variáveis forem iguais, execute uma ação (exemplo: desbloquear uma fase ou ativar um efeito).

\sphinxAtStartPar
\sphinxstylestrong{Desafio 3: Combinação de condições e variáveis}
Crie um cenário onde múltiplas condições devem ser atendidas para uma ação ocorrer.
Por exemplo, se variável A for maior que 10 \sphinxstylestrong{e} variável B for igual a 5, então o personagem ganha um bônus.
\subsubsection*{Bibliografia}

\sphinxAtStartPar
SCIRRA, L. (2022). Manual Construct.
\sphinxurl{https://www.construct.net/en/make-games/manuals/construct-3}

\sphinxAtStartPar
Rabin, S. (n.d.). Introdução ao desenvolvimento de games (Cengage Le, Vol. 1).
\sphinxurl{https://www.amazon.com.br}/Introdução\sphinxhyphen{}desenvolvimento\sphinxhyphen{}games\sphinxhyphen{}Entendendo\sphinxhyphen{}universo/dp/852211143X

\sphinxAtStartPar
Rita, S. (2017). Treinamento em Lógica de Programação (R. Sandra (ed.); Universo d).
\sphinxurl{https://br.99ebooks.net/download/treinamento-em-logica-de-programacao/}

\sphinxstepscope


\chapter{Construct: Introdução e navegação}
\label{\detokenize{modelo_aula_06:construct-introducao-e-navegacao}}\label{\detokenize{modelo_aula_06::doc}}
\sphinxAtStartPar
\sphinxstylestrong{Tempo de Aula:} 1h30min
\subsubsection*{Recursos Didáticos}

\sphinxAtStartPar
Projetor, quadro branco, computadores.
\subsubsection*{Metodologia}

\sphinxAtStartPar
Expositiva, aplicação de exercícios / solução de problemas.
\subsubsection*{Conteúdo Programático}
\begin{itemize}
\item {} 
\sphinxAtStartPar
Navegação na plataforma.

\item {} 
\sphinxAtStartPar
Criação de um novo projeto

\item {} 
\sphinxAtStartPar
Testar jogos da plataforma

\end{itemize}
\subsubsection*{Descrição}

\sphinxAtStartPar
A navegação é um ponto crítico que deve ser sempre relembrado, pois muitos alunos se
perdem dentro da área da Engine, por falta de entender esse conceito de navegação.
Criação de personagem e do cenário no chão. Adicionar os comportamentos .
No decorrer dessa aula deve\sphinxhyphen{}se deixar os alunos explorarem para desenhar o próprio
personagem e adicionar comportamentos.
\subsubsection*{Avaliação}

\sphinxAtStartPar
Execução dos exercícios em sala de aula.
\subsubsection*{Observação}

\sphinxAtStartPar
Pode ser feito uma análise para verificar se compensar ter todos os arquivos em uma
pasta para o aluno, assim ele trabalha diretamente com a ferramenta sem se preocupar em
baixar as imagens ou caçalas na internet. Isso acaba levando muito tempo de uma aula 30
\% na escolha e pesquisa. Sendo assim deixaremos uma pasta com personagens GIF e jpg
Cenários e etc.
\subsubsection*{Bibliografia}

\sphinxAtStartPar
SCIRRA, L. (2022). Manual Construct.
\sphinxurl{https://www.construct.net/en/make-games/manuals/construct-3}

\sphinxAtStartPar
Rabin, S. (n.d.). Introdução ao desenvolvimento de games (Cengage Le, Vol. 1).
\sphinxurl{https://www.amazon.com.br}/Introdução\sphinxhyphen{}desenvolvimento\sphinxhyphen{}games\sphinxhyphen{}Entendendo\sphinxhyphen{}universo/dp/852211143X

\sphinxAtStartPar
Rita, S. (2017). Treinamento em Lógica de Programação (R. Sandra (ed.); Universo d).
\sphinxurl{https://br.99ebooks.net/download/treinamento-em-logica-de-programacao/}

\sphinxstepscope


\chapter{Construct: Animação de personagens}
\label{\detokenize{modelo_aula_07:construct-animacao-de-personagens}}\label{\detokenize{modelo_aula_07::doc}}
\sphinxAtStartPar
\sphinxstylestrong{Tempo de Aula:} 1h30min
\subsubsection*{Recursos Didáticos}

\sphinxAtStartPar
Projetor, quadro branco, computadores.
\subsubsection*{Metodologia}

\sphinxAtStartPar
Expositiva, aplicação de exercícios / solução de problemas.
\subsubsection*{Conteúdo Programático}

\sphinxAtStartPar
Compreender a área de edição de sprites e animações.
Conhecer a lista e configuração de animação.
Criação de Frames dentro da Ferramenta.
Zonas de Colisão no jogo ( Habilitar no Layout )
Pontos Âncoras da imagem( pontos de imagem )
\subsubsection*{Descrição}

\sphinxAtStartPar
Iniciar um projeto do zero, buscando através da repetição fazer o aluno decorar os
caminhos, criar personagem simples no início, e animá\sphinxhyphen{}lo a partir dos frames.
É importante entender as zonas de colisão, pois com base nelas o personagem irá interagir
com o ambiente, o chão ou outros elementos.
Nessa aula é importante frisar as animações, e configurações.
\subsubsection*{Avaliação}

\sphinxAtStartPar
Execução dos exercícios em sala de aula.
\subsubsection*{Exercícios de Desafio}

\sphinxAtStartPar
\sphinxstylestrong{Desafio 1: Criar uma animação personalizada}
Utilize a ferramenta de edição de sprites para criar uma animação com pelo menos 5 frames diferentes.
Configure a velocidade da animação para que o movimento fique natural e fluido.

\sphinxAtStartPar
\sphinxstylestrong{Desafio 2: Ajustar zonas de colisão}
No seu personagem animado, habilite e ajuste as zonas de colisão para que elas estejam alinhadas corretamente com a forma do sprite em cada frame.
Teste a interação do personagem com o chão ou outros objetos para garantir que as colisões funcionem adequadamente.

\sphinxAtStartPar
\sphinxstylestrong{Desafio 3: Manipular pontos âncoras da imagem}
Experimente alterar os pontos âncoras do sprite para entender como eles afetam a posição e o comportamento da animação.
Crie um pequeno projeto onde o personagem gira ou se move com base nos pontos âncoras definidos.
\subsubsection*{Bibliografia}

\sphinxAtStartPar
SCIRRA, L. (2022). Manual Construct.
\sphinxurl{https://www.construct.net/en/make-games/manuals/construct-3}

\sphinxAtStartPar
Rabin, S. (n.d.). Introdução ao desenvolvimento de games (Cengage Le, Vol. 1).
\sphinxurl{https://www.amazon.com.br}/Introdução\sphinxhyphen{}desenvolvimento\sphinxhyphen{}games\sphinxhyphen{}Entendendo\sphinxhyphen{}universo/dp/852211143X

\sphinxAtStartPar
Rita, S. (2017). Treinamento em Lógica de Programação (R. Sandra (ed.); Universo d).
\sphinxurl{https://br.99ebooks.net/download/treinamento-em-logica-de-programacao/}

\sphinxstepscope


\chapter{Construct: Animação de Personagem II}
\label{\detokenize{modelo_aula_08:construct-animacao-de-personagem-ii}}\label{\detokenize{modelo_aula_08::doc}}
\sphinxAtStartPar
\sphinxstylestrong{Tempo de Aula:} 1h30min
\subsubsection*{Recursos Didáticos}

\sphinxAtStartPar
Projetor, quadro branco, computadores.
\subsubsection*{Metodologia}

\sphinxAtStartPar
Expositiva, aplicação de exercícios / solução de problemas.
\subsubsection*{Conteúdo Programático}

\sphinxAtStartPar
O aluno aprenderá o que é uma folha de Eventos e como utilizar as condições para
animar o personagem. Especificamente animar evento de movimentos do
personagem inserir entradas no layout como teclado ou mouse.
Espelhamento do Personagem.
\subsubsection*{Descrição}

\sphinxAtStartPar
O aluno deverá criar personagem, o chão do cenário e criar animações do
personagem, correndo, parado e golpes.
Cada Lista da animação deverá ser ativada caso alguma condição seja satisfeita.Por
exemplo: Seta para direita foi apertado, ou barra de espaço.
Nesse primeiro contato com eventos, ele deverá entender que se trata das mesmas
condições que existem no Scratch, basta lembrar os alunos.
\subsubsection*{Avaliação}

\sphinxAtStartPar
Execução dos exercícios em sala de aula.
\subsubsection*{Exercícios de Desafio}

\sphinxAtStartPar
\sphinxstylestrong{Desafio 1: Animação com múltiplos estados}
Crie animações distintas para o personagem: parado, correndo e atacando.
Configure a folha de Eventos para que o personagem alterne entre essas animações conforme as teclas pressionadas (exemplo: seta para direita para correr, barra de espaço para ataque).

\sphinxAtStartPar
\sphinxstylestrong{Desafio 2: Controle por teclado e mouse}
Implemente entradas para que o personagem possa ser controlado tanto pelo teclado (movimentos e ataques) quanto pelo mouse (por exemplo, clique para atacar).
Certifique\sphinxhyphen{}se de que as animações corretas sejam ativadas para cada ação.

\sphinxAtStartPar
\sphinxstylestrong{Desafio 3: Espelhamento do personagem}
Configure o personagem para que, ao andar para a esquerda, a imagem seja espelhada horizontalmente, dando a sensação de que ele está virado para o lado correto.
Teste a transição entre espelhamento e animações sem causar interrupções visuais.
\subsubsection*{Bibliografia}

\sphinxAtStartPar
SCIRRA, L. (2022). Manual Construct.
\sphinxurl{https://www.construct.net/en/make-games/manuals/construct-3}

\sphinxAtStartPar
Rabin, S. (n.d.). Introdução ao desenvolvimento de games (Cengage Le, Vol. 1).
\sphinxurl{https://www.amazon.com.br}/Introdução\sphinxhyphen{}desenvolvimento\sphinxhyphen{}games\sphinxhyphen{}Entendendo\sphinxhyphen{}universo/dp/852211143X

\sphinxAtStartPar
Rita, S. (2017). Treinamento em Lógica de Programação (R. Sandra (ed.); Universo d).
\sphinxurl{https://br.99ebooks.net/download/treinamento-em-logica-de-programacao/}

\sphinxstepscope


\chapter{Construct: Tile Maps}
\label{\detokenize{modelo_aula_09:construct-tile-maps}}\label{\detokenize{modelo_aula_09::doc}}
\sphinxAtStartPar
\sphinxstylestrong{Tempo de Aula:} 1h30min
\subsubsection*{Recursos Didáticos}

\sphinxAtStartPar
Projetor, quadro branco, computadores.
\subsubsection*{Metodologia}

\sphinxAtStartPar
Expositiva, aplicação de exercícios / solução de problemas.
\subsubsection*{Conteúdo Programático}

\sphinxAtStartPar
Criar cenários baseados em tilemaps prontos.(Fornecidos pela equipe)
Trabalhar mudança de cenário.
\subsubsection*{Descrição}

\sphinxAtStartPar
Ensinar ao aluno a utilizar o mosaico para criação de cenários, produzindo pelo
menos dois níveis em um jogo no estilo plataforma, um portal deve levar o aluno a outro
nível.
\subsubsection*{Avaliação}

\sphinxAtStartPar
Execução dos exercícios em sala de aula.
Modalidade: presencial.
\subsubsection*{Exercícios}

\sphinxAtStartPar
\sphinxstylestrong{Exercício 1: Criar um cenário básico com tilemap}
Utilize os tilemaps fornecidos para montar um cenário inicial para seu jogo.
Organize os tiles para formar plataformas e obstáculos básicos.

\sphinxAtStartPar
\sphinxstylestrong{Exercício 2: Criar múltiplos níveis}
Com base no cenário criado, construa pelo menos dois níveis diferentes usando tilemaps.
Cada nível deve ter características distintas, como layout ou dificuldade.

\sphinxAtStartPar
\sphinxstylestrong{Exercício 3: Implementar troca de cenário via portal}
Adicione um portal no primeiro nível que, ao ser ativado pelo personagem, leva o jogador para o segundo nível.
Configure os eventos para que a troca de cenário ocorra de forma suave e funcional.
\subsubsection*{Bibliografia}

\sphinxAtStartPar
SCIRRA, L. (2022). Manual Construct.
\sphinxurl{https://www.construct.net/en/make-games/manuals/construct-3}

\sphinxAtStartPar
Rabin, S. (n.d.). Introdução ao desenvolvimento de games (Cengage Le, Vol. 1).
\sphinxurl{https://www.amazon.com.br}/Introdução\sphinxhyphen{}desenvolvimento\sphinxhyphen{}games\sphinxhyphen{}Entendendo\sphinxhyphen{}universo/dp/852211143X

\sphinxAtStartPar
Rita, S. (2017). Treinamento em Lógica de Programação (R. Sandra (ed.); Universo d).
\sphinxurl{https://br.99ebooks.net/download/treinamento-em-logica-de-programacao/}

\sphinxstepscope


\chapter{Construct: Elementos de jogo \sphinxhyphen{} Barra de vida, Moedas (status)}
\label{\detokenize{modelo_aula_10:construct-elementos-de-jogo-barra-de-vida-moedas-status}}\label{\detokenize{modelo_aula_10::doc}}
\sphinxAtStartPar
\sphinxstylestrong{Tempo de Aula:} 1h30min
\subsubsection*{Recursos Didáticos}

\sphinxAtStartPar
Projetor, quadro branco, computadores.
\subsubsection*{Metodologia}

\sphinxAtStartPar
Expositiva, aplicação de exercícios / solução de problemas.
\subsubsection*{Conteúdo Programático}

\sphinxAtStartPar
Criar barra de vida manualmente
Criar um indicador de moedas.
Conceito de Camadas
\subsubsection*{Descrição}

\sphinxAtStartPar
Ensinar o aluno a Construir uma barra de vida, que demonstre a quantidade de vida
que o personagem possui, lembrando que é necessário criar variáveis globais que
armazenam a quantidade de vida.
Ativar a camada Paralax, 0,0 para que o personagem ao se movimentar seja seguido pela
quantidade de vida.
\subsubsection*{Avaliação}

\sphinxAtStartPar
Execução dos exercícios em sala de aula.
\subsubsection*{Bibliografia}

\sphinxAtStartPar
SCIRRA, L. (2022). Manual Construct.
\sphinxurl{https://www.construct.net/en/make-games/manuals/construct-3}

\sphinxAtStartPar
Rabin, S. (n.d.). Introdução ao desenvolvimento de games (Cengage Le, Vol. 1).
\sphinxurl{https://www.amazon.com.br}/Introdução\sphinxhyphen{}desenvolvimento\sphinxhyphen{}games\sphinxhyphen{}Entendendo\sphinxhyphen{}universo/dp/852211143X

\sphinxAtStartPar
Rita, S. (2017). Treinamento em Lógica de Programação (R. Sandra (ed.); Universo d).
\sphinxurl{https://br.99ebooks.net/download/treinamento-em-logica-de-programacao/}

\sphinxstepscope


\chapter{Construct: Aula de Revisão}
\label{\detokenize{modelo_aula_11:construct-aula-de-revisao}}\label{\detokenize{modelo_aula_11::doc}}
\sphinxAtStartPar
\sphinxstylestrong{Tempo de Aula:} 1h30min
\subsubsection*{Recursos Didáticos}

\sphinxAtStartPar
Projetor, quadro branco, computadores.
\subsubsection*{Metodologia}

\sphinxAtStartPar
Expositiva, aplicação de exercícios / solução de problemas.
\subsubsection*{Conteúdo Programático}

\sphinxAtStartPar
REFORÇAR CONTEÚDOS DAS AULAS PASSADAS.
\subsubsection*{Descrição}

\sphinxAtStartPar
Reforço dos conteúdos passados, visando levar o aluno a relembrar os principais
conceitos das aulas.
\subsubsection*{Avaliação}

\sphinxAtStartPar
Execução dos exercícios em sala de aula.
\subsubsection*{Bibliografia}

\sphinxAtStartPar
SCIRRA, L. (2022). Manual Construct.
\sphinxurl{https://www.construct.net/en/make-games/manuals/construct-3}

\sphinxAtStartPar
Rabin, S. (n.d.). Introdução ao desenvolvimento de games (Cengage Le, Vol. 1).
\sphinxurl{https://www.amazon.com.br}/Introdução\sphinxhyphen{}desenvolvimento\sphinxhyphen{}games\sphinxhyphen{}Entendendo\sphinxhyphen{}universo/dp/852211143X

\sphinxAtStartPar
Rita, S. (2017). Treinamento em Lógica de Programação (R. Sandra (ed.); Universo d).
\sphinxurl{https://br.99ebooks.net/download/treinamento-em-logica-de-programacao/}

\sphinxstepscope


\chapter{Construct: Hierarquia e Camadas}
\label{\detokenize{modelo_aula_12:construct-hierarquia-e-camadas}}\label{\detokenize{modelo_aula_12::doc}}
\sphinxAtStartPar
\sphinxstylestrong{Tempo de Aula:} 1h30min
\subsubsection*{Recursos Didáticos}

\sphinxAtStartPar
Projetor, quadro branco, computadores.
\subsubsection*{Metodologia}

\sphinxAtStartPar
Expositiva, aplicação de exercícios / solução de problemas.
\subsubsection*{Conteúdo Programático}

\sphinxAtStartPar
HIERARQUIA DE SPRITES
\subsubsection*{Descrição}

\sphinxAtStartPar
Entender o que é hierarquia, anexar filhos aos pais de cada elemento.
\subsubsection*{Avaliação}

\sphinxAtStartPar
Execução dos exercícios em sala de aula.
1 \sphinxhyphen{} O personagem principal deve possuir um objeto que será anexado a si mesmo.
\subsubsection*{Exercícios}

\sphinxAtStartPar
\sphinxstylestrong{Exercício 1: Criar hierarquia de objetos}
Crie um personagem principal e anexe a ele um objeto filho, como uma espada ou escudo.
Faça com que o objeto filho se mova e gire junto com o personagem principal.

\sphinxAtStartPar
\sphinxstylestrong{Exercício 2: Manipular camadas para profundidade}
Organize diferentes objetos em camadas distintas (por exemplo, personagem em uma camada, cenário em outra).
Teste a visibilidade e sobreposição para garantir que os objetos apareçam na ordem correta durante o jogo.
\subsubsection*{Bibliografia}

\sphinxAtStartPar
SCIRRA, L. (2022). Manual Construct.
\sphinxurl{https://www.construct.net/en/make-games/manuals/construct-3}

\sphinxAtStartPar
Rabin, S. (n.d.). Introdução ao desenvolvimento de games (Cengage Le, Vol. 1).
\sphinxurl{https://www.amazon.com.br}/Introdução\sphinxhyphen{}desenvolvimento\sphinxhyphen{}games\sphinxhyphen{}Entendendo\sphinxhyphen{}universo/dp/852211143X

\sphinxAtStartPar
Rita, S. (2017). Treinamento em Lógica de Programação (R. Sandra (ed.); Universo d).
\sphinxurl{https://br.99ebooks.net/download/treinamento-em-logica-de-programacao/}



\renewcommand{\indexname}{Índice}
\printindex
\end{document}